\section{General Conference} % *GENERALCONFERENCE

\subsection{October 2015}

\subsubsection{Things I'd like to work on or change}

\begin{itemize}
\item Larry Lawrence's challenge to ask what is holding us back.
\item Save money each week (Durrant?). See his talk for more info.
\item Ponderize one verse of scripture per week. Choose a verse and place it where you'll see it each day. Then think of it each day and its phrases throughout the week. Do it weekly for six months or longer. I'll do it for a while.
\end{itemize}

\subsubsection{Larry Lawrence (Mackenzie's dad?)}

Ask in prayer, very soon, what is holding me back from progressing? Then quietly listen to the Spirit.

Direction more important than speed.

Cites a Kimball quote about piercing the veil? Check this out.

\subsection{April 2015}

\subsubsection{Things I'd like to work on or change}

\begin{itemize}
\item I want to start listening to BYU and gospel devotional addresses again. I will think of some times that are appropriate for me to do this. First up: Oaks' talk ``Our strengths can become our downfall.''
\item Re-read Bednar.
\item Re-read Christofferson.
\item Re-read, with Alicia, the couple of talks on Sunday morning that dealt with doubt, questions, and people who have left the church.
\end{itemize}

\subsubsection{Sister Burton (relief society president)}

She spoke about the way that husbands and wives treat each other. To me it's a good reminder to speak to Alicia in a kind way. I was praying the other night about this---specifically, about how I need to be more quick to forgive and more quick to soften my heart when we are ending fights. This is something that I am going to pray about and work on.

\subsubsection{Elder Oaks}

Elder Oaks spoke about the parable of the sower, and he went through Christ's explanation of the parable and then drew parallels to the way we live in the world today and the way we are converted. He talked about the \textit{keyhole} view of the gospel---this is when you are focused on a single doctrine, principle, or deficiency in the church or in a leader. You're too focused on that one thing. That reminds me of the metaphor that I think Oaks himself has used in the past, of the person who plays a single gospel note on the piano over and over and over.

He talks about that in a devotional address called ``Our strengths can become our downfall.'' Looks like he gets the metaphor from Boyd Packer. Here is the quote:

\begin{quote}
My first example concerns Satan's efforts to corrupt a person who has an unusual commitment to one particular doctrine or commandment of the gospel of Jesus Christ. This could be an unusual talent for family history work, an extraordinary commitment to constitutional government, a special gift in the acquisition of knowledge, or any other special talent or commitment.

In a memorable message given at the 1971 October conference, Elder Boyd K. Packer likened the fulness of the gospel to a piano keyboard. He reminded us that a person could be ``attracted by a single key,'' such as a doctrine they want to hear ``played over and over again.'' He explained:

\textit{Some members of the Church who should know better pick out a hobby key or two and tap them incessantly, to the irritation of those around them. They can dull their own spiritual sensitivities. They lose track that there is a fulness of the gospel, . . . [which they reject] in preference to a favorite note. This becomes exaggerated and distorted, leading them away into apostasy. [Boyd K. Packer, \textit{Teach Ye Diligently} (Salt Lake City: Bookcraft, 1975), p. 44]}

We could say of such persons, as the Lord said of the members of the Shaker sect in a revelation given in 1831, ``Behold, I say unto you, that they desire to know the truth in part, but not all'' (D\&C 49:2). And so, I say, beware of a hobby key. If you tap one key to the exclusion or serious detriment of the full harmony of the gospel keyboard, Satan can use your strength to bring you down.
\end{quote}

\subsubsection{Elder Bednar}

It was about fear; re-read this.

\subsubsection{Elder Christofferson}

Not quite sure what it was about, Addy was having a really hard time. Re-read.

\subsubsection{Elder Holland}

In this talk he suggests that there's no reason to celebrate Easter or Christmas unless you think that Adam and Eve actually existed and actually fell. Is that true? I'm still not sure that you really need to think that, though I'm willing to and it seems more interesting to think that they actually did exist, and then try to work out what that would mean.

Also, a few years ago he gave a talk in which he said something about how Satan really does exist, as well.

\subsection{October 2014}

Things to look up:

\begin{enumerate}
\item Quote given by Christofferson from Brigham YOung, I think, about how each member needs a testimony when following the prophet.
\item Re-read Christofferson with Alicia.
\item Re-read Uchtdorf, Saturday morning, with Alicia.
\item Re-read Oaks, Saturday afternoon, with Alicia
\item Re-read Anderson, Saturday afternoon---it's a pretty direct answer to my question about how to help others. He gives some good suggestions but I think his most important point is to fix your own mask before you help others. That is, get your own faith and your own house in order before you go out to try and help other people in their doubts. Unless you're ready to go yourself, you're not going to be able to help anyone.
\item re-read callister sat after. lots about teaching.
\item print out and re-read J\"{o}rg's talk.
\item Re-read Uchtdorf's priesthood talk, this directly addresses my question about being a gospel doctrine teacher.
\end{enumerate}

\subsection{April 2014}

I didn't really come into this conference with any questions. The main point I want to resolve is: what graduate program should I choose? I have been praying about it for a few days now, and feel that Duke is the right choice, and that's the conclusion that I have presented to the Lord. I will be fasting tonight (Saturday, April 5), to look for confirmation of this decision. I'm mainly interested in thoughts about prioritizing family life and living the gospel in the world; perhaps other comments will also make the difference in helping me decide. From everything that I have read, looked at, and learned in person, Duke seems like the right place. 

\paragraph*{Saturday Morning}

\subparagraph*{Thomas Monson}

How does he not get tired of the cultural performances before a temple dedication? He talks about one in Arizona. Does Arizona really have a culture or a history, besides Hispanic or Mexican history?

\subparagraph*{Jeffrey Holland}

Alternative to a religious life is a life left desolate. 

People want gods; they just want gods that are easy to live with. They want gods that don't demand anything from them. They want it easy. Christ did not preach comfortable doctrine. It is not comfortable to be a disciple of Christ. It takes hard work. If it isn't hard, if it doesn't require sacrifice, then I'm probably not doing it right.

``I am more certain that the keys of the priesthood have been restored...than I am certain that I am standing before you and you are in the congregation in front of me.'' Pretty remarkable thing to say. This is not the exact quote, but contains the idea.

\subparagraph*{Rasband}

\subparagraph*{Carlos H. Amado}

\subparagraph*{Linda Reeves}

Pornography. Family plan with standards and boundaries to protect home from porn. I should have a relationship good enough with Addy and the other kids that they could talk to me if they saw something, or were having a problem. I don't really know how to develop this type of relationship with them.

\subparagraph*{Neil Andersen}

\subparagraph*{Henry Eyring}

\subsection{April 2013}

I had no questions to keep track of during this conference.

\paragraph*{Saturday Morning}

\subparagraph*{President Monson} New temples: Cedar city, Rio de Janeiro

\subparagraph*{Boyd Packer} Be careful of the tolerance trap--tolerance for some acts does not reduce the spiritual conqeuences of those acts when they go against what God says. We have to stay worthy of the Holy Ghost to receive its promptings and pure intelligence (``stay in condition to respond to the Holy Ghost''). We can know the things we need to know instantly.

\subparagraph*{Dean M. Davies}

\subparagraph*{Elaine S. Dalton} The part of young women matters because they matter (to God). Their responsibilities may seem mundane but they are important. Story of her mother, who did great things for her family but was never recognized. Quotes family proclamation --role of father, role of mother. Difference in the use of ``primary''? Discrepancy? Return to virtue. Privilege of being a woman, wife, and a mother.

Think about double standards in evaluating these talks --if a man says something weird, we just say, he's a man. If he says something good for women then we praise him. If a woman says something weird, we say she's a shill; if she says something good, we don't praise him, but we expect it.

\subparagraph*{Craig Cardon} The Savior forgives sin on earth and at the judgment. Recurring human weakness is not beyond the power of the Savior to save. The Lord doesn't hold back forgiveness according to the seriousness of the sin. He allows for improvement over time, not immediate perfection. When we repent and seek forgiveness, he forgives. The Savior wants to forgive.

\subparagraph*{M. Russell Ballard} The Lord used the priesthood to endure the atonement? Very interesting. Check this quote. Priesthood is different from priesthood authority. The authorization or ordination is given by the laying on of hands. The power of the priesthood comes only when those who exercise it are worthy and acting in accordance with God's will. Men have the unique responsibility to administer the priesthood. But they are not the priesthood. Men and women have different but equally valued roles. Just as woman cannot conceive a child without a man, so a man cannot fully exercise the power of the priesthood to establish an eternal family without a woman. In other words, in the eternal perspective, both the procreative power and priesthood power are shared by husband and wife. And both husband and wife should strive to follow our Heavenly Father... (get more). Puts emphasis on the new videos released by the church on strengthening the family and the church through the priesthood.

\subparagraph*{Henry Eyring} As we are faithful to the promises we have made with God, we will feel love for Him. As I obey, the Savior will draw close to me and bless my family with what is best. That seems like a direct enough answer to questions about the future. Just keep being obedient and see what's in store for us.

\paragraph*{Saturday afternoon}

\subparagraph*{Richard Scott} Regardless of your circumstances, you can make the Lord the center of your life. Apply sound gospel principles in life; don't take shortcuts. Little things lead to big things, and little neglects can turn into big problems. Simple and consistent habits lead to full and bountiful blessings. Satan's fate is decided; he knows he is lost, but wants to take as many with him as he can. MEMORIZE SCRIPTURES. I can do this. Just need to incorporate it again into flash cards. When we obey the commandments of the lord and serve his children unselfishly, the natural consequence is power from God. Power to do more than we can do by ourselves. Our insights, our talents, are expanded, because we receive strength and power from the Lord. His power is a fundamental component to establishing a home filled with peace. As you center your home on the Savior, it naturally becomes a refuge not only to your own family, but also to friends who live in more difficult circumstances. They will be drawn to the serenity they feel there. Welcome such friends into your home. They will blossom in that Christ-centered environment. One of the greatest blessings we can offer to the world is the power of a Christ-centered home.

\subparagraph*{Quentin Cook} Difference between universal and individual peace. They have called the term ``agency'' by the name ``moral agency'' twice in this conference so far. That phrase appears in the scriptures but I don't think it's very common elsewhere, at least I don't hear it very much.

It is as individuals and families that we achieve the peace that is the promise of righteousness. It's an abiding, deep happiness and spiritual contentment. I looked up the phrase ``moral agency'' in the general conference corpus --it is far more common beginning in the 1960's, and the most common of all in the 90's and the 00's. Repentance and living righteously is essential for peace. ``Spiritual but not religious'' --that's a good first step, but in the church we're fellowshipped and nourished, taught the word of god, and can partake of priesthood ordinances that prepare us to return to God's presence. These ordinances bring peace because they are covenants with the Lord. Because of the atonement, righteous living will be rewarded with peace.

\subparagraph*{Stanley Ellis} First article of faith: nothing more direct. We can learn from our heavenly parents, and do other things because of them. We aren't spiritual orphans. Helping the poor is great, but in helping we can hurt. The way is self-reliance; we shouldn't do for others what they can do for themselves. JS taught correct principles, and the members governed themselves. Only once in 16 years did he get a call asking, which ward needs help? Which ward needs a good family? We should do all we do in his way.

\subparagraph*{John B. Dickson} Peter received the guidance to take the gospel to others. The expansion of the church was made known to its presiding officers and then to the general membership. The concept was new to the saints, that it could go to others. Peter got the revelation as the senior apostle. In the time of the restoration of the gopsel, JS restored the church and the first presidency moved the church forward. Things have moved forward according to the Lord's calendar and his time; then, in 1978, came the relveation on the priesthood, and it went to Kimball, the prophet. This time, ALL the people in the world could partake of the gospel. The Lord has established a pattern for this, that is true.

\subparagraph*{David A. Bednar} Why is the law of chastity so important? We have experiences with a body that we couldn't have otherwise. Stuff on men and women; need to get these quotes. Every impulse of the natural man can be overcome through the atonement. Love increases through righteous constraint, and decreases through indulgence.

\subparagraph*{Russel Nelson} Addresses ward mission leaders directly.

\paragraph*{Priesthood session}

\subparagraph*{Robert Hales} The whole armor of God is our only protection.

Plan to have a goal-setting session --and plan to do it regularly. Sometimes we're the lightning rod and we must take the heat for living the gospel --but we are never alone.

\subparagraph*{Callister} Still strange to me that I've already served a mission --what was that training and retraining? Did I get that? Did I do it right? I can barely remember now. If I am a youth leader, let the kids take charge.

\subparagraph*{Beck (?)} Virtually everyone I know could benefit in some way from my ministering.

\subparagraph*{Dieter Uchtdorf} ``Moral agency'' again --4 times I think.

\subparagraph*{Eyring}

\subparagraph*{Monsen} Future missionaries: are you willing to work? ready to serve? It requires selfless sacrifice, long hours, etc.

Pretty much this whole meeting was about the Aaronic priesthood and missionary work (or preparing for it).

Plan life with a purpose!

\paragraph*{Sunday morning session}

\subparagraph*{Dieter Uchtdorf} Darkness exists, but we don't have to dwell there. Light also exists, and we can choose to dwell in it. God's light is real. It is available to all. It gives life to all things. It has the power to soften the sting of the deepest wound. It can be a healing balm for the loneliness and sickness of our souls. It can enlighten the deepest valleys of sorrow. It can lead us through the dark night. It is the spirit of Jesus Christ. But spiritual light doesn't come to those who wait for someone to flip the switch; it takes an act of faith. How do we begin? Start where we are. The very first steps we take to the light bring immediate good into our lives. Walk in the light. He does not wish to break our spirits; he wishes us to rise up.

\subparagraph*{Neil Andersen} Put my old missionary tag where I can see it. I have a contribution to make. In everyday conversations, I can add to the faith of others. Hold up Jesus Christ. Names and faces will com into my mind as I pray to know to whom to speak. The words will come to me too. Opportunities will be open to me. Faith will overcome doubt. The Lord will bless me with my very own miracles.

\subparagraph*{Rosemary Wixom} Turn everything off sometimes to connect with children, to give them full attention.

\subparagraph*{Whitney Clayton} I am married. How should I do? First, in the happiest marriages, both the husband and wife consider their relationship to be a pearl beyond price, a treasure of infinite worth. They build a marriage that will prosper for eternity. Their path is divinely ordained. No other relationship of any kind can bring as much joy, do as much good, or bring as much personal refinement. Second, faith. Successful eternal marriages are built on the foundation of faith in the Lord and adherence to his teachings. Couples who have made their marriage priceless practice the patterns of faith. Weekly attendance at meetings, prayer, pay tithing, read scriptures. Keep the commandments. Third, repentance. Happy marriages rely on the gift of repentance. It is an essential element in every good marital relationship. Spouses who regularly conduct honest self examinations and promptly take steps to improve experience a healing balm in their marriages. Repentance restores harmony and peace. Humility is the essence of repentance. Selfless, not selfish. Doesn't demand its own way or speak with moral superiority; it listens, it recognizes hat no one can change someone else. With faith we can undergo our own mighty change of heart. This causes us to treat our spouses with meekness. Humility means that both husbands and wives seek to bless, help, and lift each other, putting the other first in every decision. Fourth, respect. Husbands and wives treat each other as equal partners. Husbands shouldn't dominate. Equal voice and vote. Focus first on the home. Marriages are based on cooperation. Dinner time and following are objects of best effort. Forego personal entertainment to help with household duties. Fifth, transparency. No secrets about relevant matters in marriage. Husbands and wives make decisions about finances together. Fiercely loyal to each other. No secret internet experiences. They never do or say anything that approaches the appearance of impropriety. Sixth, love. Leave behind single life. Marriage is first priority.-Marriages are gift from God to us; the quality of our marriage is a gift from us to God.

\subparagraph*{Tom Perry} Secularism is beginning to be the norm. Many of its beliefs and practices are in direct conflict with the things instituted by the Lord for the blessing of his children. ``Obedience to law is liberty.'' Satan is driven to succeed; do not underestimate this. There are moral absolutes: sin will always be sin, disobedience will always bring consequences. I should look this quote up, very interesting. The moral absolutes aren't about particulars; they're about a species of universals, in some general statements. I like this a lot. We should not pick and choose which commandments to keep; we should acknowledge all of them.

\subparagraph*{President Monson} When we keep the commandments, we receive truth and light. That's it. We just do. Knowledge of truth and the questions of our existence come to us as we are obedient to the commandments of God. Remember, to obey is better than to sacrifice. Obedience comes first. Obedience is a source of strength and knowledge. Knowledge of what? Experience? This is a very interesting type of knowledge. Seems different from propositional knowledge...this needs to be elaborated. It's important I think. The great test of this life is obedience.

\paragraph*{Sunday afternoon}

\subparagraph*{Jeffrey Holland} An urgent, emphatic desire is enough to begin. What matters is not the size of our faith; the integrity of our actions toward the faith that we do have, and to the knowledge that we do have, is what matters. Don't let questions stand in the way of the faith that we do have working a miracle. Don't lead with a declaration of doubt; that is not superior to a declaration of a small amount of faith. What we know always trumps what we don't know. In this world, everyone is to walk by faith. When you see imperfection, the limitation is not because of the lack of divinity in the work. You don't have to apologize for ``only believing.'' This is a fallacy in the church --perhaps because of the influence of ``knowing'' in science. Belief matters. Consider the articles of faith: they start with ``we believe,'' not ``we know.'' Honestly acknowledge your questions and concerns, but first fan the flames of your faith, for everything is possible to those who believe.

\subparagraph*{Dallin Oaks}

\subparagraph*{Christoffel}

\subsection{October 2012}

I had three questions for these sessions. Here they are:

\begin{enumerate}
\item How can I learn, write, and do philosophy with the Holy Ghost as my guide?
\item How can I better approach situations like Tandi's, where someone close to me is asking serious questions about the church? What is my role in these situations? What should I do?
\item How can I gain power through temple worship?
\end{enumerate}

I'm not sure that these notes will be as formally presented as those from last conference; I'm not so much interested in the presentation as I am in whether or not I return to the notes and make them a part of my study in the next six months.

\paragraph*{Saturday Morning}

\subparagraph*{Cook} I need to be able to answer Alma's ``Can you feel so now?'' question affirmatively. Experience a change of heart daily --seek the Spirit daily. Email Tandi about Cook's quote --magnifying and inventing weaknesses of early church leaders. I can be erotic and sexual with Alicia without using degrading language or speaking in ways that make me feel uncomfortable later. Tell her how beautiful she is in as many different ways as I can, and show her I want her.

Look up Isaiah scripture --run and not be weary, walk and not faint.

\subparagraph*{Dibb} Purpose of knowledge: life and godliness. Knowledge is not an end in itself.

\subparagraph*{???} Tandi: live to have the Holy Ghost. I need to be guided by the Spirit. Remember that I have the gift of the Holy Ghost, available to me all the time through the baptismal covenant.

\subparagraph*{Bowrein (?)}

\subparagraph*{Nelson} Get quote about truth from the laboratory/revelation.

\subparagraph*{Uchtdorf (REREAD)} Men in particular regret not spending time with family. Remember this --print it out. Get the quote and print it.

The Son of God lives purposefully each day. Don't just pretend to spend time with people.

Take notes again with questions here!

Question 1: Remember that my family and service are more important.

\paragraph*{Saturday Afternoon}

\subparagraph*{Perry} Email Nate --use quotes about how it's the parents' responsibility to teach kids. It's theirs to teach to think too.

\subparagraph*{Ballard} Minds and hearts must be engaged during the week.

Our testimony must transcend our mind and burrow into our heart. Get quote.

Question 1: assimilate doctrines of the gospel deep into my soul. See promised blessings.

In my morning prayer, ask Heavenly Father to guide me to someone I can serve. See and get this quote.

See ``discern'' quote --like my blessings.

\subparagraph*{Larry Echo Hawk} I need to stand proud at Tufts and in the world as a Mormon. That is what I am.

\subparagraph*{Other guy (?)} Question 2: follow spiritual promptings wherever they lead. Ask, then act. Get ``sophisticated neutrals'' quote.

\subparagraph*{Temple talk (REREAD)} Temple standard. Temple is the house of the Lord, not mine. Things are done his way there. How can I do things in the Lord's way there? Get ``unholy temple'' quote. Can't hide sins from the Lord --prepare to visit by repenting. What do I need to repent of? in order to stand in holy places.

\subparagraph*{Anderson (REREAD)} Get quotes on internet church stuff. Read about what he says of the law of chastity.

\subparagraph*{Oaks}

\paragraph*{Saturday Evening Priesthood}

Write in journal!

Write down commitment to never sell a book that has anything to do with the church. If the ALR is ever finished, it will be free.

\subparagraph*{Christofferson} Question 1: Act in faith. Remember my blessings --seek the kingdom first, and the things I need will be supplied for me.

\subparagraph*{Stevenson}

\subparagraph*{Perkins} Question 3: In temple worship, we may receive a fullness of the Holy Ghost. Get that exact quote.

\subparagraph*{Uchtdorf}

\subparagraph*{Eyring} Pray to know what gifts I have, how to use them, and how to help others discover their gifts also.

\subparagraph*{Monsen} Some don't have a testimony. Share mine to help them grow theirs.

\paragraph*{Sunday Morning}

\subparagraph*{Eyring (REREAD)} In rereading, look for barriers to revelation and feeling God's love. Find Keaton/Chris Riley.

\subparagraph*{Packer}

\subparagraph*{Linda Burton} First observe, then serve.

\subparagraph*{Walter F. Gonzalez (REREAD)} Question 2: Remember --impressions will give me knowledge that I can't get any other way. Remember Alma 26:22 --need to continue seeking the Spirit to tell them what the Lord wants them to know.

Gather quotes about knowledge from this conference.

Question 1 --knowledge from heaven helps in other subjects --GET QUOTE

\subparagraph*{Holland} If I love Christ, why am I still on the same shores, doing the same old things? Move onto greater things.

Judgment day Q: Did you love me? Did we at least understand that commandment?

Crowning characteristic of love: loyalty.

\subparagraph*{Monson} Take inventory of life, look for blessings, large and small, that I have received.

Much of the 2 Nephi 2:25 joy comes from recognizing that we can communicate in prayer and receive answers.

Write in journal.

Thank you notes --be more aware of people and how I can show them that they are important to me.

\paragraph*{Sunday Afternoon}

\subparagraph*{Hales} Denying ourselves of ungodly behavior is first step of repentance.

\subparagraph*{Scott} Any temple work done in the temple is good --but need to do my own. Question about temple work: my own ancestor's work.

Put family history on planning list.

\subparagraph*{Osgelthorpe}

\subparagraph*{Marcus Nash}

\subparagraph*{Daniel Johnson}

\subparagraph*{Don Clarke} Benefit from meetings: come with problems, be humble and willing to listen, and have a desire to help those around us.

\subparagraph*{Bednar (REREAD)} Receiving testimony/being converted --conversion requires more than knowing.

Definition of testimony --get it. True conversion brings a change of beliefs (?).

Answer to questions 1, 2, and 3: personal righteousness.

``Setting aside weapons requires more than believing and knowing.''

\subsection{April 2012}

The main question that I wish to take into this conference is this one: what specific things can I change and do in order to have the Spirit in greater measure? A couple things will help me reflect on this question, I think. First, are worldly concerns overly important in my life? Second, are my home and my life places where the Spirit can dwell? Third, how can I invoke the blessing of the Lord being close to me before I begin my career? These are just a few thoughts to guide my listening. The most important question is, what specific things can I do to have the Spirit?

This time, instead of making one gigantic list of things to do and going with that, at the end of conference I'll compile a list of things I can do and then go through and prioritize it. I'll set the items into four or five groups--the most important group, the next most important group, etc. Then I'll work on doing or implementing into my life those items in the first group, and when I've been successful with that, I'll do the second group, and so on. This way I won't be overwhelmed and I will do the most important things first.

So, the driving question is this, followed by a few points to guide my thinking:

What specific things can I change and do in order to have the Spirit in greater measure?

Are worldly concerns overly important in my life? Which ones?

Are my home and life (my mind) places where the Spirit can dwell?

How can I invoke the blessing of the Lord being close to me before my career starts?

These will guide me. Lately I have been thinking often about the importance of personal revelation and the Spirit's company throughout the day. This is the main reason that I chose this question and have had this problem in mind. I want to live by every word that proceeds forth from the mouth of God. In order to do that, I need to have the Spirit more, and I need to be more inclined to obey and follow its promptings and direction. The generic recipe for having the Spirit more is to be more righteous, but as I go into this conference, I am looking for specific things that I can change and do to be more righteous and have the Spirit more.

Items in red are those which answer the above questions and are things I can do in order to have the Spirit more.

Items in green are things that I need to schedule time to do or think about.

Items in blue are passages or entire talks that I need to read again.

\paragraph*{Saturday Morning}

\subparagraph*{President Monson} I must oppose evil wherever it is. How can I do this without being a zealot or bigot? I don't want to ruin relationships--I think that would do more harm than good. I need to build constructive relationships that generate understanding, but I also must have a clear idea of what I believe and why and in what ways these beliefs may depart from typical or traditional LDS beliefs.

\subparagraph*{President Packer} Read scriptures with Alicia every day--pray before and after--pay attention and be thoughtful and open to discussion and to the Spirit. This will bring the Spirit into our home--I need to open my heart in these situations! The time with my family is sacred--I can make this time more special by changing my attitude.

\subparagraph*{Cheryl Esplin}

\subparagraph*{Donald Hallstrom} The church and the gospel are not the same--the gospel is the plan of salvation of the children of God, through Jesus Christ's sacrifice.

The church was established by Christ, restored through Joseph Smith (by Joseph Smith's hand). Christ both was and is the head of his church.

But what does it mean to say that the church and the gospel are not the same? Is it that the church ``administers''? What does that mean?

There is something more than the church--it is the gospel.

Activity in the church is not the ultimate goal--it is possible to be active in the church and inactive in the gospel. The church is an outward thing, but the things of the gospel are usually less visible and difficult to measure, and they are of much greater importance. These are things like, how much faith do we have? How repentant are we? How important are our covenants?

The purpose of the church is to help us live the gospel.

We can make the gospel the foundation in our life by doing the following things:
\begin{enumerate}
\item Deepening our understanding of deity;
\item Focusing on ordinances and covenants, and by establishing the discipline to live them faithfully;
\item Uniting the gospel and the church--if we focus on the gospel, church will be more meaningful. The Lord wants members of his church to be converted to and active in the gospel.
\end{enumerate}

Become active in the gospel.

\subparagraph*{Paul Koelliker} How can I be less distracted by worldly things?

Awakening the desire to know more about the gospel--it's the quest of each of us--this desire will help me focus on what I need to do to avoid distractions.

[he read some scriptures here--find them and review]

Awakening desire --> Seek pattern --> Learn doctrine of Christ

\subparagraph*{Oaks} The atonement doesn't end sacrifice in the gospel plan; see 3 Nephi 9:20. Our sacrifice is of our time and our selfish priorities.

Our sacrifice is on a day to day basis in our ordinary lives.

The sacrifice requires effort--it will be difficult, like being humble with my family.

``It's true, isn't it? Then what else matters?'' But I'm not expected to give all my time [like a monk or nun], I'm expected to make daily sacrifices in my time and selfish priorities (and feelings/reactions!). That is how I live the gospel at work, school, and everywhere.

[reread the part about lectures on faith]

\subparagraph*{Eyring} ``Give me these challenges--mountains.'' I can achieve my goals. I can achieve them.

The solid basis for the foundation of faith is personal integrity, choosing the right consistently. Getting older doesn't make the foundation stronger--serving God does. It does take time, but that isn't sufficient.

\paragraph*{Saturday Afternoon Session}

\subparagraph*{Holland} Don't feel envy when someone else is fortunate--it diminishes me not at all. Envy is the enemy, not good fortune.

``Be kind and grateful that God is kind--it is a happy way to live.'' Just don't hold on to mortal failings.

I doubt the fact that I was born into a great family is due to anything that I have deserved or earned. It is just a blessing, an occurrence motivated by God but not understood by me. Other people are no less than I am. I just got lucky, in a way--and so of him unto whom much is given, much is required.

\subparagraph*{Hales} Keep the Sabbath day holy--this invites the Spirit into my mind and home, and this spirit will allow me to ponder on the Atonement. Not just the facts of it, but how it relates to me.

The purpose of spiritual and temporal self-reliance is to move us to higher ground so that we may lift others.

\subparagraph*{Baxter}

\subparagraph*{Ulisses Soares} Let go of all negative emotions--how can I do this?

\subparagraph*{Cook} Avoid being overly judgmental (or judgmental at all, in most cases) about behavior that is foolish or unwise but not sinful.

The distance between God and me narrows and spirituality returns when I immerse myself in the scriptures.

\subparagraph*{Scott (this entire talk contains things I can do and focus on)} Revelation is crisp, clear communication from the Holy Ghost.

Inspiration is a prompting or set of promptings toward a worthy objective.

I need to learn to ask. This is how revelation is received. Fast and pray to find and understand scriptures--it's a cyclical process: read, ponder, pray, ponder, pray, receive impressions)--a good way to learn from the scriptures.

Yielding to emotions like anger, hurt, and defensiveness drives away the Holy Ghost. Eliminate these emotions from myself.

Loud laughter offends the Spirit, but a sense of humor is good.

Exaggeration offends the Spirit--careful, quiet speech is good.

Health, exercise, good eating, and good sleeping increase one's ability to receive spiritual communication.

Strive to capture content in (dreams?)

Record impressions of the Spirit--this shows God that the communication is sacred, and it enhances recall. Be careful with sharing them.

Sanctify yourself (in the scriptures) = keep my commandments.

One whom the Lord trusts has access to inspiration.

Haughtiness, pride, and conceit are stony ground--humility is the fertile soil. Arrogance or a desire to be seen and heard of others or a person who lets his emotions control his actions is not a condition for the Spirit.

Indicators of good inspiration (inspiration from God): peace in heart and a quiet, warm feeling.

The disposition to do right brings peace and happiness.

The conditions for receiving personal revelation are: obedience, trust in the plan of happiness, avoidance of anything contrary to the plan.

Revelation is not trivial--help comes in response to obedience, faith, and our wise use of agency.

\paragraph*{Saturday Priesthood Session}

\subparagraph*{Bednar} The priesthood has no selfish component--it's all about serving others. When we have it, we are things that act, we are not things that are acted upon. We act with initiative.

I must work diligently in my new calling--it's a new chance!

I need power in the priesthood.

Seriously, this is the priesthood--the power of God, given to man. What am I doing, not magnifying my calling? The reward is exaltation!

I am the priesthood leader in our home--this is my responsibility. How can I do this better?

I should get my line of authority from Dad.

\subparagraph*{Edgley} Pray about the passages in my blessings that talk about discerning needs and serving those who are less active.

\subparagraph*{Ochoa} Actively examine the media and entertainment in which I participate. Remove those things that do not invite the Spirit.

Now is the time to take my potential and make it real.

\subparagraph*{Uchtdorf} With my calling, Tufts, and other things, I just need to step out into the dark and have faith, doing my best.

Inspired leaders help others to see the why of the gospel, and help others to get to work.

That something is good is not sufficient for it to require our time--stay centered in discipleship--self-discipline is needed to choose those actions which will be of most benefit--read scriptures, pray diligently to be guided by the Spirit.

Sermons that do not lead to action are worthless. Actually living the doctrine is what matters.

Kneel and ask--beg--supplicate--entreat--for help and guidance. Then follow it!

God and Christ live. They are real. They are there.

\subparagraph*{Eyring} Same as Bednar--this is the priesthood, the sealing power of God. The sealing power of God!

He explains the keys of Elias--about ten minutes in. 

Discusses the Holy Spirit of Promise, about twelve minutes in (quotes someone else).

More stuff about fathers and their responsibility.

\begin{enumerate}
\item Get and preserve a testimony that they keys are here, on earth, with Pres. Monson
\item Love my wife--put her needs above mine in all things--my love will grow
\item Get the help of everyone to love one another
\item ??? [Read the end of section 121---the Spirit will be my constant companion]
\end{enumerate}

\subparagraph*{Monson} Check out the John Taylor quote (early in the talk)--so what is the deal with mapping the gospel and science to each other?

Don't murmur. Don't complain.

``When God commands, and a man obeys, that man is always in the right.''

When unselfish service replaces selfish concerns, miracles happen.

``You've written to him 17 times and he's never responded?'' ``So? Maybe this will be the month.''

Look at this man's life. It's amazing. And though I can't have the same experiences, I can work just as hard. There's nothing that will stop me.

\paragraph*{Sunday Morning Session}

\subparagraph*{Uchtdorf} We assume that we have all the information we need to make appropriate judgments about others--but we rarely do. We assign dark motives to others in order to justify our bitterness, coldness, or anger, and our lack of forgiveness.

STOP IT. Replace judgmental thoughts and feelings with a heart full of love for God and his children.

Get Keaton's email to apologize. And Chris Riley.

Allow God's love to govern our mind and actions--darkness of animosity and envy will eventually fade as we let God's love in.

Stop gossiping.

\subparagraph*{Nelson} Gratitude--how can I be more aware of God's providence, and express love?

\subparagraph*{Rasband}

\subparagraph*{Beck}

\subparagraph*{Christofferson} Establishing doctrine is a matter of divine revelation.

How the Savior communications knowledge--reread (about halfway).

Not every statement made by a church leader should be considered doctrine.

\subparagraph*{Monson} Start memorizing.

Divine attributes: to think, to reason, to achieve.

\paragraph*{Sunday Afternoon Session}

\subparagraph*{Perry}

\subparagraph*{Ballard} Promptings of Holy Ghost bring us back to the path, and to the Atonement.

We are prosperous because we have strong families--the families are the cause, not the effect.

The most important cause of our life is families--everything in our lives and society will improve if we can strengthen families.

Life is better and much happier when hearts turn to families.

Put all outside concerns in subjection to what happens in families.

Read the proclamation to the family.

As I seek to live the doctrine of Christ, the Spirit will teach me what to do.

\subparagraph*{Haleck}

\subparagraph*{Wilson} ``They that are wise...are not deceived (those that take the Holy Spirit as a guide?)''. A scripture?

\subparagraph*{Evans} Write letters to Brock about Jesus Christ.

Missionary experiences must be current--to be fulfilled, we must continue to naturally and normally share the gospel.

Pray to know who I can help in the quorum, and promise to obey. Then OBEY.

\subparagraph*{Paul Pieper} Our determination depends on how we treasure that which we receive from above.

The idea of attributing impressions to reason or intuition--impressions are subtle, natural, spirit-to-spirit communication.

Record my impressions. How?

That which is sacred to God becomes sacred to us through agency. Sacred/secular is an issue of priority, not exclusivity.

\subparagraph*{Andersen} Is it possible for others to consider me a Christian? What thinks Christ of me?

This life is a long migration toward a more celestial life. We're not at our best every day, but daily duty, done at our best, is the path.

Start keeping track of contacts, friends, people--write them down, track them, remember them, contact them.

\subparagraph*{Monson} Make plans to review these notes!

\subsection{October 2011}

\begin{enumerate}
\item How can I better receive and develop the gift of discernment?

Inquire of the Lord diligently--ask in faith, keep the commandments, desire, don't harden my heart. I won't receive it every time nor receive every detail. Ways in which the revelation can come: enlightening mind, peace to mind, mind and heart (Holy Ghost), burning bosom, feel it is right, fill soul with joy. Seek this gift in conjunction with scripture study. (Barbara Thompson)

Act. I don't need a planning meeting. I know what is right to do, I don't always need the gift of discernment. (Jose L. Alonso)

Look around and use your own judgment too--who needs help? Who looks like they have a particular need? (Carl B. Cook)

The spirit of love and hope will accompany me as I share my beliefs with others. (L. Tom Perry)

Look for instances of the gift of discernment in journals. (W. Christopher Waddell)

Communicate regularly through prayer. (President Monson)

I need to be worthy to receive it, then follow it. (President Monson)

Personal study on the gift of discernment

New Testament scriptures:

Matthew 16:3, \textcolor{red}{\textgreek{diakr'inw}}, ``distinguish, settle, decide.''

1 Corinthians 2:14, \textcolor{red}{\textgreek{anakr'inw}} ``examine closely, inquire into, examine (as to qualifications), examine persons concerned in a suit.'' The word is in the passive voice in this verse.

Hebrews 4:12, \textcolor{red}{\textgreek{kritik'os}} ``able to discern, for judging.''

All the other uses of ``discern$\ast$'' are forms of \textcolor{red}{\textgreek{diakr'inw}}, I think. I haven't double checked this but I think it's right.

The most relevant verse here is 1 Corinthians 2:14. Hebrews 5:14 is also good. I need to pay close attention to these verses.

I think it's important to note that the phrase ``gift of discernment'' does not appear in the scriptures; actually, the word ``discernment'' doesn't appear either. What we get is ``discerning,'' and when we're talking about spiritual gifts, we mean 1 Corinthians 12:10 and DC 46:23. The Greek word in 1 C 12:10 is the noun form of \textcolor{red}{\textgreek{diakr'inw}}, so we might say that it means ``a discerning'' or ``a distinguishing,'' but we could probably say that it just mean ``discernment'' in general.

But the question is, are these two verses (1 C 12:10 and DC 46:23) talking about the gift of discernment? In the church we talk about the gift of discernment almost as though it were a kind of mind-reading, an ability to understand the thoughts and the needs of others in order to help them come unto Christ. I don't think that this is the same as the ``discerning of spirits,'' at least in the understanding of past members of the church--I think that gift is meant to be a protection to individuals which allows them to see evil designs and intentions in the world and avoid them. They are not quite the same. I think we ought to understand the gift of discernment as an instance of the spirit of revelation, given in special circumstances and sometimes associated with the stewardship that one has over a person or group of people. But here are some quotes to take a look at.

Stephen Richards, General Conference 1950:

\begin{quotation}
First, I mention the gift of discernment, embodying the power to discriminate, which has been spoken of in our hearing before particularly as between right and wrong. I believe that this gift when highly developed arises largely out of an acute sensitivity to impressions -spiritual impressions, if you will -to read under the surface as it were, to detect hidden evil, and more importantly to find the good that may be concealed. The highest type of discernment is that which perceives in others and uncovers for them their better natures, the good inherent within them. It's the gift every missionary needs when he takes the gospel to the people of the world. He must make an appraisal of every personality whom he meets. He must be able to discern the hidden spark that may be lighted for truth. The gift of discernment will save him from mistakes and embarrassment, and it will never fail to inspire confidence in the one who is rightly appraised.

The gift of discernment is essential to the leadership of the Church. I never ordain a bishop or set apart a president of a stake without invoking upon him this divine blessing, that he may read the lives and hearts of his people and call forth the best within them. The gift and power of discernment in this world of contention between the forces of good and the power of evil is essential equipment for every son and daughter of God. There could be no such mass dissensions as endanger the security of the world, if its populations possessed this great gift in larger degree. People are generally so gullible one is sometimes led to wonder whether the great Lincoln was right, after all, in the conclusion of his memorable statement, ``You can't fool all the people all the time.'' One does feel at times, however, a sense of pity and sympathy for some of the peoples of the world whose education, information, and exposure to higher ideals and exalted concepts have been so arbitrarily and ruthlessly restricted.

There is a class of people now grown sizable in the world who should possess this great gift in large degree. They know how the gift is attained. They have been educated in its spiritual foundations. They have been blessed with the counsels which foster it. They know how to order their lives to procure it. You know who they are, my brethren and sisters. Every member in the restored Church of Christ could have this gift if he willed to do so. He could not be deceived with the sophistries of the world. He could not be led astray by pseudo-prophets and subversive cults. Even the inexperienced would recognize false teachings, in a measure at least. With this gift they would be able to detect something of the disloyal, rebellious, and sinister influences which not infrequently prompt those who seemingly take pride in the destruction of youthful faith and loyalties. Discerning parents will do well to guard their children against such influences and such personalities and teachings before irreparable damage is done. The true gift of discernment is often premonitory. A sense of danger should be heeded to be of value. We give thanks for a set of providential circumstances which avert an accident. We ought to be grateful every day of our lives for this sense which keeps alive a conscience which constantly alerts us to the dangers inherent in wrongdoers and sin (General Conference, 1950).
\end{quotation}

S. Dilworth Young, Seventy, Conference 1952:

\begin{quotation}
I sustain him [David O. McKay] as a prophet. I do not believe that everything that he says is prophetic, but I earnestly pray always that I may have the gift of discernment, so that when he does speak prophetically I shall be able to recognize it and follow the instruction given.
\end{quotation}

Marion G. Romney, Conference 1956:

\begin{quotation}
In conclusion, I again call attention to the statement of the Prophet Joseph Smith already quoted, that ``A man must have the discerning of spirits before he can drag into daylight this hellish influence and unfold it unto the world in all its soul-destroying, diabolical, and horrid colors;'' (TPJS 211), for after all, the things of God can be understood only by the spirit of God. (See 1 Cor. 2:11.) The gift of ``discernment of spirits'' is the sure solution to this knotty problem. Seek after this gift, brethren and sisters, and after its kindred gifts -knowledge, wisdom, and ``to know the diversities of operations whether they be of God,'' -and not after sensational and miraculous signs and wonders.
\end{quotation}

George Q. Cannon, JD 26:

\begin{quotation}
Let the Teachers visit under the influence of the Spirit of God and the gift of discernment, and where they find those that are living in opposition to, or in violation of the laws of God, let them, by the Spirit of God, which will rest upon them, teach and warn that household, and thus take steps to purify the Church.
\end{quotation}

\item How can I gain a greater desire (causually efficacious motivation) for eternal life?

Read the Book of Mormon with humility and faith. (Richard G. Scott) % comment

Pray to receive the fruits of the spirit when I minister to others. (Jose L. Alonso)

Read mission journals and compile and type important passages. (Jose L. Alonso)

Do things to see the glory of the earth and universe--walks, movies, telescope, etc. All things in space and time are created for ordinary humans like us. Reflect on these worlds without end. (President Uchtdorf)

Make plans to remember and thank and notice others, like a little book where I keep a record of the people that I meet. I can put something in the planning about thinking on the places where I will be going that day and reviewing the book to see if I know anyone there who I might run into. (President Uchtdorf)

Become a master manager of my time. (Ian S. Ardern)

President Monson looks up to the important things--he has an attitude of ``can do, will do.'' He execises great faith and looks up to the glory of life and the incomprehensible glory of the future. (Carl B. Cook)

Receive strength through the Atonement and remove burden of sins and impurities through repentance--get rid of that stuff so it doesn't cloud my life--Keaton, Chris Riley, what else? (Carl B. Cook)

Repentance is the cause for true celebration. (D. Todd Christofferson)

The spirit of love and hope will accompany me as I share my beliefs with others. (L. Tom Perry)

Don't forget the converts and members in Chile--reach out to them in greater measure. (W. Christopher Waddell)

Consider the lessons learned and gifts obtained in my mission. How can I apply them more in post-mission life? (W. Christopher Waddell)

Remember the face without pain, fear, or guilt. (W. Christopher Waddell)

Prepare others to return to God's presence with me. (President Eyring)

Make the decision to stand alone, if necessary, in defense of the church and gospel. (President Monson)

Look for a calling card to call Mirzhi. (President Monson)

Read the Book of Mormon--focus on the Atonement--hope through it for eternal life. (Tad R. Callister)

Organize the Book of Mormon reading for our Sunday school class. This will require me to be organized, and I will have to make it a priority every week. If it's going to happen, I have to stay on top of it.

Work on loving Alicia. The mission can't be had again but I can create that type of life and relationship with my wife. (Elaine S. Dalton)

Don't compromise the desire for eternal life by doing things that I know will hurt my spirituality. (President Monson)

Desire--ask and it shall be given; in other words I need to make sure to keep praying for it. (Dallin H. Oaks)

\item How can I be more consistent in my personal spirituality?

Memorize scriptures during the day--on Kindle? (Richard G. Scott)

Make a plan to read the Old Testament. (Richard G. Scott)

Read the New Testament every day with Alicia. (Richard G. Scott)

Those who read the Book of Mormon regularly have special blessings: added measure of the Spirit of the Lord, a greater resolve to obey His commandments, and a stronger testimony of the divinity of the Son of God. (Richard G. Scott)

Ultimately, I must make the decision and just do it in the best way I can think of. (Ian S. Ardern)

Get ready to work effectively every day. (W. Christopher Waddell)

Remember the face without pain, fear, or guilt. (W. Christopher Waddell)

Form a habit of pushing through fatigue and fear. Spiritual staying power comes from pushing past the point where others would rest. Fundamental. President Hinckley one asked President Eyring, ``Why are you sleeping?'' President Monson always asks, ``Am I up to date on my work?'' Discipline myself to work harder than I thought I could. Get into better physical shape to have more energy to work. I can push myself beyond my limits of endurance and develop that as a habit. (President Eyring)

Remember that it is all a process of becoming. (President Eyring)

\end{enumerate}

All Notes

\paragraph*{Saturday Morning}

\subparagraph*{Richard G. Scott} We do not need to be concerned with the validity of concepts in the standard works: ``Because scriptures are generated from inspired communication through the Holy Ghost, they are pure truth. We need not be concerned about the validity of concepts contained in the standard works since the Holy Ghost has been the instrument which has motivated and inspired those individuals who have recorded the scriptures.'' Scriptures are ``pure truth,'' since they are given by the Holy Ghost. Scriptures give renewed confidence and power in one's ability to deal with the challenges of life. How can I give the Christmas gift to my dad of that scripture?

\subparagraph*{Barbara Thompson} Is the gift of discernment coextensive with the gift of revelation/gift of receiving revelation?

\subparagraph*{L. Whitney Clayton}

\subparagraph*{President Monson} General temple patron assistance fund

\subparagraph*{Jose L. Alonso}

\subparagraph*{Boyd K. Packer} Gender established in the premortal world? He quotes the family proclamation and adds his own comment: ``Gender is an essential characteristic [and was established in that premortal existence]...'' Voice of the spirit comes as a feeling, rather than a sound. ``Feeling'' is the most appropriate word to describe spiritual communication.

\subparagraph*{President Uchtdorf} Don't give into extremes--that man is untouchable or that man is worthless. Neither is true. The Lord doesn't care where we work--only that we are doing our best with our heart inclined to him (cf. My patriarchal blessing). Listen, remember, care to do good and bless others. The pain and flaws I experience now will not be the norm in eternity.

\paragraph*{Saturday Afternoon}

\subparagraph*{David A. Bednar} Godhead, Book of Mormon, salvation/exaltation. Joseph Smith: ``The spirit, power, and calling of Elijah is, that ye have power to hold the key of the...fullness of the Melchizedek priesthood...; and to...obtain...all the ordinances belonging to the kingdom of God.'' Make real plans to do family history. My heart will turn to the fathers. The link to the fathers in my patriarchal blessing will be more meaningful, and I will be protected against the intensifying influence of the adversary.

\subparagraph*{Neil L. Anderson} We need to get serious about deciding when to have kids after April 15.

\subparagraph*{Ian S. Ardern} Set goals and put them out to see. Make detailed plans! This will be the only way to do all the things I want to do. Eliminate procrastination. Ask: What do I do that gives me no value? Eliminate those things. Prioritize with a pyramid, maybe. The Lord can help me make decisions about priorities. Being busy should be being productive. Remember: the tools of the internet are means to achieve good, productive, meaningful goals. Kick things off the saddle and work to achieve productive, good goals. Do this--analyze the way I spend time. Wisely eliminate distractions. Stop saying I'll do something tomorrow, next week, whatever.

\subparagraph*{Carl B. Cook}

\subparagraph*{LeGrand R. Cook Jr.} Invite others to serve in the church--invite them to repent and come back.

\subparagraph*{D. Todd Christofferson} Justification = forgiveness of sins, sanctification = purification from the effects of sin. Repentance is striving to change--God's grace is a complement and reward for our greatest effort. Daily effort may be required, don't worry. Pray for time and opportunities to overcome. Repentance is not box-checking--making lists of steps might be helpful to some but it is really not a mechanical process. The Holy Ghost brings relief and forgiveness. Any pain in repentance will always be far less than the suffering for unresolved sins to satisfy justice. The morning of joy upon accomplishing goals is worth any suffering which is required.

\subparagraph*{L. Tom Perry} Be bold in the declaration of Jesus Christ. Be righteous examples to others. Speak up about the church. I should try to keep a respectful and Christlike tone in these conversations. Communicate the message honestly and directly. Share these beliefs!

\paragraph*{Priesthood session}

\subparagraph*{Jeffrey R. Holland}

\begin{enumerate}
\item Satan is real--the personification of evil
\item Satan is eternally opposed to the love of God, the atonement, the work of peace and salvation.
\end{enumerate}

He had a great influence over Joseph in his first prayer. ``Sign on, speak up''--we are at war. We need missionaries to open their mouths. To play, we must be morally clean. Personal worthiness is the greatest and most impenetrable defense. ``Satan is a real being set on destroying you.'' I need to be ready to serve another mission--get myself back into shape.

\subparagraph*{Keith B. McMullin} The priesthood is not duty-free. Power in the priesthood comes from personal righteousness and exercised faith. Live to enjoy the Holy Ghost--in order to receive power. From my righteousness, God will endow my priesthood with power.

\subparagraph*{W. Cristopher Waddell} The Lord knows which mission president, companions, members, and investigators we need to become the greatest we can be. Talk about with Alicia: ``Missionary work is taxing, difficult, etc.''

\subparagraph*{President Uchtdorf} Personal, family, and church priorities should spring from the two ``great'' commandments. Our spiritual progress is inseparably bound to the temporal service we give to others. Important story about President Kimball's stake getting flooded. Handouts do not develop self-reliance! This is my responsibility. The gospel does not work by magic. Study the principles and doctrines, study the manuals, use provident living.org. Then, apply what I learn to those who are in my stewardship. If I follow this pattern, I will receive specific guidance on how to do it.

\subparagraph*{President Eyring} The challenge will come when the good memories fade--be an example of the believers.

\subparagraph*{President Monson} There is no doubt about what is moral. There is no doubt about what is right and wrong. It is not personal, it is not relative, it is not cultural, it is not a matter of taste. We can know it and we need to stay in that position and say it. He [President Monson] stood alone at times--he made the decision to do it. In the long run, you can't have your cake and eat it too. Failures in this area are flaws of thinking--the philosophy we carry around has a great effect! ``If we know the Book of Mormon is true, it follows that Joseph was a prophet, he saw God and Christ, etc...''

\paragraph*{Sunday Morning}

\subparagraph*{President Eyring} Becoming charitable--feelings of charity become part of our very nature, in time. As we serve others, God lets us feel of his love, and then feelings of charity become part of our nature, in time. Righteousness and the spirit inside me over time will kill the desire for things which are contrary to the gospel of Jesus Christ.

\subparagraph*{Elder Hales} Wait upon the Lord = hope, trust, anticipate. Have patience in the Lord--the plan takes a lifetime for a reason. Be kind to one another, kinder to ourselves.

\subparagraph*{Tad R. Callister} Bible/Book of Mormon line example. Book of Mormon is either true or not.

\subparagraph*{Elaine S. Dalton} The most important thing I can do for my daughter: love Alicia. I need to eliminate anything from my life which is unclean or inconsistent with the priesthood.

\subparagraph*{Elder Ballard} The name of Christ, the son of God--we take it upon ourselves. Cf. The original name of the priesthood--something is going on there. Use the full name of the church.

\subparagraph*{President Monson} Morality is dying, but we don't despair. Our code of conduct is not negotiable. Nothing brings more joy or peace than the spirit which comes from obedience. ``It is essential to reject anything which does not conform to our standards.''

\paragraph*{Sunday afternoon}

\subparagraph*{Elder Nelson}

\subparagraph*{Elder Oaks} What do I believe about Christ? What does it cause me to do?

\subparagraph*{Matthew Richardson} In teaching, let the Holy Ghost do it--teach by the spirit. I may need to change the way I teach to align it with the Holy Ghost's.

\begin{enumerate}
\item The Holy Ghost teaches in a very personal way. Line upon line, precept on precept, here and there--you teach people, not lessons. Focus on those things which people need to know and do.
\item The Holy Ghost teaches us by inviting and encouraging us to act.
\end{enumerate}

\subparagraph*{Kazuhiko Yamashita}

\subparagraph*{Randall K. Bennett} Must choose eternal life.

\subparagraph*{J. Devon Cornish} The things which are important to us are important to God--story of the letter to Cristopher.

\subparagraph*{Elder Cook} Trials come but we need to see the whole perspective.

\subparagraph*{President Monson} Pray for the general authorities.

Things I need to do:

\begin{enumerate}
\item Review the three questions above, and their answers, regularly.
\item Make plans to type up and read mission journals, looking in particular for instances of the gift of discernment.
\item Do things to see the glories of the universe with my wife.
\item Set up my telescope.
\item Set up a system (a little book?) to notice, remember, and thank others.
\item Apologize to Keaton, Chris Riley, and any others I think of.
\item Reach out to converts and members in Chile.
\item Reflect on the gifts I received in my mission. How can I apply them more in post-mission life?
\item Look for a calling card to call Mirzhi.
\item Memorize scriptures during the day--on Kindle? In the mornings? Not sure what the best plan for this is.
\item Make plans to read the Old Testament in 2012.
\item Develop the habit of pushing through fatigue and fear.
\item Figure out how I can give the Christmas gift to my dad of that scripture.
\item Make plans to actually do family history on Sunday.
\item Analyze the ways I spend my time. What can I eliminate? What is of no value? What things are of the most value, and how can I do more of them?
\item Invite the people I visit to repent and come back to activity.
\end{enumerate}
% -----------------------------------------------------------
\section{Meeting Notes}
% -----------------------------------------------------------

\section{Sacrament Meetings}

\subsection{January 2013}

Paul Carlisle spoke today. My thoughts when he was speaking: religion doesn't exist to answer questions, at least not ultimately. It exists to give us the possibility of a spiritual life, and that takes work. God isn't up there just throwing out answers to difficult questions. He's up there to be a participant in our spiritual life with us; if we do not make an effort, then we don't have a spiritual life. That's just the way it works. If people don't like that because it's different from science or they think that it shouldn't be that way for one reason or another, then too bad. The fact that things work a certain way in the world does not say anything about how religion does or should work. It works a certain way, whether we like it or not.

\subsection{August 2013}

\subsubsection{August 25}

Sarah Briggs spoke today, along with Rob Wessman and Christian Cahoon. Sarah spoke about hope and optimism. I'm realizing more and more that optimism is a choice \textit{and} a responsibility: if we really believe that good things are going to happen, and we understand that something of our part is required to make them happen, then we will get to work. So optimism carries a responsibility to work toward the good things that we want. So I guess the way to tell a real optimist is to check who's working toward the things they really want.

Austin Boaz gave the elders quorum lesson. We talked about making your spiritual life the center of your life; it will have various components, one of which is family, and another of which can be a professional life, I think, but the spiritual life should be the center.

\subsection{September 2013}

\subsubsection{September 1}

During priesthood meeting today, President Bowen (of the temple presidency) said that most people are in ruts--they do the same old things and don't reflect on what they are doing in order to change. He said, ``If I were a young father, I would want to teach my children how to learn from their experience--how to learn experientially. That requires reflection and categorization, observing and listening.''

Paul Carlisle talked about how our goal should be to expand our mind \textit{and} heart as we learn. He used the example of Enoch, whose mind and heart expand when the Lord shows him things. If we are only expanding our mind, that's the learning of men, and that isn't what we want. Spiritual learning comes by study and by faith and obedience. That is how we gain intelligence.

So maybe we can use the study/faith distinction to drop some other concepts into the two sides.

\begin{table}
\centering
 \begin{tabular}{p{4cm}|p{4cm}}
 Study \& Faith and Obedience \\
 Mind (Enoch's mind expands with new knowledge) \& Heart (Enoch''s spiritual commitment deepened) \\
 Learned (when they are learned, they think they are wise) \& Wisdom (when they are learned, they think they are wise) \\
 \end{tabular}
\end{table}

You can probably fill in a lot more than that, as well. I like this distinction but I really like being able to put some more meat in it by adding other concepts.

Tom Eager reiterated that knowledge must come by study and by faith. You can't just \textit{study} (mind only) whatever you want and expect to be fine. He mentioned how, six weeks ago, he had brought up Out-stitute as an example of the type of learning you shouldn't engage in, and how some people in the stake were upset that he had even mentioned it at all.

\subsection{January 2014}

\subsubsection{January 5}

A general heuristic or principle for thinking about the gospel: we should question any ``doctrine,'' idea, belief, etc., which grants a special status or privilege to a person or group of people only in virtue of (i) their country of birth, (ii) family of origin, (iii) region of birth, (iv) era of birth, (v) religion of birth.

A scriptural example of this is in Mt 3:9---don't say, we have Abraham for our father.